\documentclass[12pt]{article}

\usepackage[utf8]{inputenc}
\usepackage{amsmath, amssymb, amsthm}
\usepackage{geometry}
\geometry{a4paper, margin=2.5cm}

\title{Măsura Dirac}
\author{}
\date{}

\theoremstyle{plain}
\newtheorem{theorem}{Teoremă}

\begin{document}

\maketitle

\section{Definiție}

Fie $(X, \mathcal{B})$ un spațiu măsurabil și fie $x_0 \in X$ un punct fixat.
\textbf{Măsura Dirac în punctul $x_0$}, notată $\delta_{x_0}$, este aplicația:



\[
\delta_{x_0} : \mathcal{B} \to \{0,1\}, \qquad
\delta_{x_0}(A) = 
\begin{cases}
1, & \text{dacă } x_0 \in A, \\
0, & \text{dacă } x_0 \notin A.
\end{cases}
\]



Aceasta este o măsură de probabilitate concentrată în întregime în punctul $x_0$.

\section{Proprietăți}

\begin{enumerate}
    \item $\delta_{x_0}$ este o măsură de probabilitate:
    

\[
        \delta_{x_0}(X) = 1.
    \]



    \item Este complet atomică, cu un singur atom:
    

\[
        \delta_{x_0}(\{x_0\}) = 1.
    \]



    \item Pentru orice funcție măsurabilă $f : X \to \mathbb{R}$,
    integrala față de măsura Dirac este:
    

\[
        \int_X f \, d\delta_{x_0} = f(x_0).
    \]



    \item Pentru orice mulțime măsurabilă $A$,
    

\[
        \delta_{x_0}(A) = \mathbf{1}_A(x_0),
    \]


    unde $\mathbf{1}_A$ este funcția indicator.
\end{enumerate}

\section{Teoremă fundamentală}

\begin{theorem}[Caracterizarea măsurii Dirac]
Fie $\mu$ o măsură de probabilitate pe $(X, \mathcal{B})$.
Atunci $\mu$ este o măsură Dirac dacă și numai dacă există un punct $x_0 \in X$ astfel încât:



\[
\mu(A) = \mathbf{1}_A(x_0), \qquad \forall A \in \mathcal{B}.
\]



În acest caz, $\mu = \delta_{x_0}$.
\end{theorem}

\section{Exemplu}

Pe spațiul $(\mathbb{R}, \mathcal{B}(\mathbb{R}))$, măsura Dirac în punctul $a \in \mathbb{R}$ este:



\[
\delta_a((-\infty, x]) =
\begin{cases}
0, & x < a, \\
1, & x \ge a.
\end{cases}
\]



Aceasta este funcția de repartiție a unei variabile aleatoare degenerate în $a$.

\end{document}
